\section{Quyết định kiến trúc hệ thống}\label{sec:architecture-decisions}
Mục này ghi nhận các quyết định kiến trúc (Architecture Decision Record -- ADR) định hình cách hệ thống TMC được xây dựng. ADR là quy tắc bắt buộc: một quyết định chỉ được thay đổi bằng một ADR mới thay thế nó. Mỗi ADR gồm bốn phần: bối cảnh, quyết định, hệ quả (lợi ích và đánh đổi) và các phương án thay thế đã được cân nhắc. Các ADR được xây dựng từ phân tích ở Mục~\ref{sec:architecture-structure} và chỉ sử dụng các yêu cầu, ràng buộc, bất biến đã được định nghĩa ở các phần trước.

\subsection{Bảng tổng hợp quyết định kiến trúc}\label{sec:adr-summary}

\begin{table}[H]
\centering
\caption{Tổng hợp quyết định kiến trúc}\label{tab:adr-summary}
\small
\renewcommand{\arraystretch}{1.3}
\begin{tabularx}{\textwidth}{|>{\centering\arraybackslash}p{1.4cm}|>{\raggedright\arraybackslash}X|>{\centering\arraybackslash}p{1.7cm}|>{\raggedright\arraybackslash}p{4.4cm}|>{\centering\arraybackslash}p{1.6cm}|}
\hline
\textbf{Mã} & \textbf{Quyết định} & \textbf{Trạng thái} & \textbf{Yêu cầu / ràng buộc liên quan} & \textbf{Mục liên quan} \\
\hline
ADR-01 & Topology triển khai kết hợp biên -- trung tâm, Gate Agent trên máy chủ biên của từng nhà máy & Accepted & NFR-PER-01, FR-17, NFR-AVL-04, NFR-REL-03, C-02, C-04, A-05 & \ref{sec:physical-step1} \\
\hline
ADR-02 & Mã QR động có chữ ký số, xác thực tại biên & Accepted & NFR-SEC-02, NFR-SEC-05, NFR-SEC-06, NFR-PER-01, R-04, R-06 & \ref{sec:quantum} \\
\hline
ADR-03 & Lưu bền vững tại cổng trước khi xác nhận và đồng bộ bằng transactional outbox & Accepted & NFR-REL-03, NFR-AVL-03, NFR-DIN-02, INV-05, INV-06, R-05 & \ref{sec:cc-views} \\
\hline
ADR-04 & Architecture style kết hợp: modular monolith tại trung tâm, microkernel và ports and adapters tại Gate Agent & Accepted & K1--K8; NFR-MOD-01 & \ref{sec:style-decision} \\
\hline
ADR-05 & Giao tiếp giữa hai quantum qua \textit{sync-contract} bất đồng bộ; giao tiếp trong Central qua interface và sự kiện trong tiến trình & Accepted & FR-17, NFR-SEC-06, NFR-DIN-02 & \ref{sec:logical-step4} \\
\hline
ADR-06 & Tích hợp cổng xoay và HR/Payroll qua adapter, có giới hạn thời gian chờ, circuit breaker và gửi lại & Accepted & C-01, C-02, NFR-REL-01, NFR-REL-02, NFR-INT-01 -- NFR-INT-04 & \ref{sec:logical-step4} \\
\hline
ADR-07 & Tổ chức mã nguồn: monorepo, phân chia theo module, kiểm thử kiến trúc tự động & Accepted & NFR-MOD-01, NFR-SCA-01 & \ref{sec:module-views} \\
\hline
ADR-08 & Kiểm soát truy cập, nhật ký kiểm toán bất biến và bảo vệ dữ liệu cá nhân & Accepted & C-03, NFR-SEC-01 -- NFR-SEC-04, NFR-AUD-01 -- NFR-AUD-05, NFR-PRI-01 -- NFR-PRI-04 & \ref{sec:physical-step4} \\
\hline
ADR-09 & Đồng bộ dữ liệu nhân viên từ HR/Payroll & Accepted & FR-01, NFR-SEC-06, NFR-INT-01 & \ref{sec:cc-views} \\
\hline
ADR-10 & Nền tảng công nghệ: Java và Spring Boot cho backend, React cho giao diện & Accepted & K6, K7; NFR-MOD-01 & \ref{sec:allocation-views} \\
\hline
\end{tabularx}
\end{table}

\subsection{Nội dung chi tiết các quyết định kiến trúc}

%------------------------------------------------------------------------------
\subsubsection{ADR-01: Topology triển khai kết hợp biên -- trung tâm}\label{adr:01}
\textbf{Trạng thái:} Accepted. \textbf{Nhóm:} cấu trúc, triển khai.

\noindent\textbf{1. Bối cảnh.} Hệ thống phục vụ hai nhà máy với khoảng 1.400 nhân viên. Giờ giao ca có khoảng 800 lượt quét trong 30 phút mỗi nhà máy, tương đương trung bình 0,44 lượt/giây và đỉnh giả định 2--5 lượt/giây mỗi nhà máy. Kiosk phải hiển thị kết quả trong 2 giây (NFR-PER-01) và cổng phải tiếp tục hoạt động khi mất kết nối với trung tâm (FR-17), trong khi đường truyền WAN giữa nhà máy và trung tâm có thể gián đoạn (R-02, A-05). Đồng thời, dữ liệu chấm công phải được tổng hợp tập trung để truyền sang HR/Payroll trước 12:00 ngày 1 (C-02).

\noindent\textbf{2. Quyết định.}
\begin{itemize}
    \item Áp dụng topology kết hợp biên -- trung tâm. Gate Agent (Gate Quantum) chạy trên \textbf{máy chủ biên của từng nhà máy} (Factory Edge Server), gồm một node chính và một node dự phòng nóng; kiosk tại trạm cổng chỉ là thiết bị đầu cuối không lưu dữ liệu.
    \item Central Quantum chạy tại trung tâm dữ liệu on-premise của TMC (A-05), trên \textbf{hai Central Application Server chạy song song} sau Reverse Proxy, cùng CSDL PostgreSQL có bản sao dự phòng.
    \item Toàn bộ đường xử lý lượt quét nằm trong mạng nhà máy; không có lời gọi đồng bộ nào qua WAN trên đường này.
\end{itemize}

\noindent\textbf{3. Hệ quả.}
\begin{itemize}
    \item \textit{Lợi ích:} cổng tiếp tục hoạt động khi mất kết nối với trung tâm; thời gian xử lý lượt quét không phụ thuộc chất lượng WAN; dữ liệu vẫn được quản lý tập trung; một kho cục bộ chung cho cả nhà máy bảo đảm quy tắc chống quẹt lặp chiều nhất quán (INV-04).
    \item \textit{Đánh đổi:} phải vận hành và cập nhật phần mềm cho các node tại nhà máy; dữ liệu giữa hai quantum chỉ nhất quán cuối cùng (độ trễ $\le$ 60 giây); thiết bị đặt tại hiện trường cần được bảo vệ (mã hóa dữ liệu, chứng chỉ thiết bị -- ADR-08). Central phải chạy hai node vì việc cấp mã QR động nằm trên đường đi của nhân viên (ADR-02).
\end{itemize}

\noindent\textbf{4. Phương án thay thế.}
\begin{itemize}
    \item \textit{Tập trung hoàn toàn:} bị loại vì mỗi lượt quét phụ thuộc WAN, vi phạm FR-17 khi mất kết nối.
    \item \textit{Phân tán hoàn toàn tại từng nhà máy:} bị loại vì dữ liệu, tài khoản và audit bị phân mảnh, khó đáp ứng C-02.
    \item \textit{Gate Agent chạy trên từng kiosk:} bị loại vì quy tắc chống quẹt lặp chiều phải đồng bộ ngang giữa các kiosk, số node cần quản lý và bảo vệ tăng theo số trạm.
\end{itemize}

%------------------------------------------------------------------------------
\subsubsection{ADR-02: Mã QR động có chữ ký số, xác thực tại biên}\label{adr:02}
\textbf{Trạng thái:} Accepted. \textbf{Nhóm:} bảo mật, hiệu năng.

\noindent\textbf{1. Bối cảnh.} Mã QR tĩnh có thể bị chụp lại, chia sẻ hoặc phát lại để ra vào trái phép hoặc chấm công hộ (R-06). Trạng thái thu hồi lưu tại trạm cổng có thể chưa kịp cập nhật (R-04), nhất là khi mất kết nối. Mặt khác, việc xác thực phải hoàn tất tại cổng, không gọi trung tâm (ADR-01), và nằm trong ngân sách 2 giây (NFR-PER-01). Mọi nhân viên đều có điện thoại để mở TMC Employee Portal (A-02).

\noindent\textbf{2. Quyết định.}
\begin{itemize}
    \item Mã QR hiển thị cho nhân viên là một token gồm phiên bản, mã credential, thời điểm phát hành, thời điểm hết hạn, một giá trị ngẫu nhiên dùng một lần (nonce) và mã định danh khóa, được ký bằng \textbf{Ed25519}. Thời hạn hiệu lực không quá 60 giây (NFR-SEC-05); Portal lấy mã mới khoảng mỗi 30 giây qua \texttt{QrTokenAPI} của C2.
    \item Khóa bí mật chỉ nằm tại Central (C2). Gate Agent chỉ nhận khóa công khai qua chiều đồng bộ xuống.
    \item G1 xác thực mã hoàn toàn bằng dữ liệu cục bộ: kiểm tra chữ ký, thời hạn (cho phép lệch đồng hồ $\pm$30 giây, các node đồng bộ thời gian qua NTP), nonce chưa được sử dụng, credential đang hiệu lực, nhân viên đang hoạt động, trạm đang bật và quy tắc chống quẹt lặp chiều.
    \item Khóa ký được thay theo chu kỳ; mã định danh khóa cho phép khóa cũ và khóa mới được chấp nhận song song trong thời gian chuyển đổi.
\end{itemize}

\noindent\textbf{3. Hệ quả.}
\begin{itemize}
    \item \textit{Lợi ích:} ảnh chụp mã QR chỉ dùng được trong thời hạn ngắn và chỉ một lần; thu hồi có hiệu lực gần như ngay lập tức kể cả khi trạm đang mất kết nối, vì Portal không cấp mã mới cho credential đã thu hồi; thiết bị tại nhà máy bị đánh cắp cũng không thể tạo mã hợp lệ; xác thực chỉ mất vài mili-giây.
    \item \textit{Đánh đổi:} việc cấp mã phụ thuộc Central, nên Central phải có hai node chạy song song (ADR-01); điện thoại của nhân viên cần kết nối mạng để lấy mã, nếu không thì bảo vệ mở cổng thủ công (FR-10); các node phải đồng bộ thời gian; cần quy trình quản lý và thay khóa.
\end{itemize}

\noindent\textbf{4. Phương án thay thế.}
\begin{itemize}
    \item \textit{Mã QR tĩnh có chữ ký (HMAC):} bị loại vì mã bị lộ dùng được vô thời hạn và khóa bí mật phải đặt tại trạm cổng.
    \item \textit{Mã TOTP sinh trên điện thoại:} không chọn vì cần ứng dụng di động lưu khóa bí mật trên thiết bị của nhân viên; có thể xem xét trong tương lai.
    \item \textit{Tra cứu trung tâm ở mỗi lượt quét:} bị loại vì phụ thuộc WAN (vi phạm ADR-01, FR-17).
\end{itemize}

%------------------------------------------------------------------------------
\subsubsection{ADR-03: Lưu bền vững tại cổng trước khi xác nhận và đồng bộ bằng transactional outbox}\label{adr:03}
\textbf{Trạng thái:} Accepted. \textbf{Nhóm:} độ tin cậy, toàn vẹn dữ liệu.

\noindent\textbf{1. Bối cảnh.} Hệ thống không được làm mất giao dịch chấm công (NFR-REL-03), không được báo mở cổng thành công khi cổng xoay chưa xác nhận (NFR-AVL-03), và không được để xảy ra trường hợp cổng đã mở nhưng không có dấu vết hoặc giao dịch bị ghi trùng (R-05). Kết nối giữa nhà máy và trung tâm có thể gián đoạn bất kỳ lúc nào.

\noindent\textbf{2. Quyết định.}
\begin{itemize}
    \item Trình tự xử lý tại Gate Agent (luồng F1): (1) xác thực; (2) ghi bền vững lượt quét ở trạng thái chờ xác nhận; (3) gọi cổng xoay; (4) chỉ khi cổng xác nhận mở mới ghi giao dịch chấm công và bản ghi outbox trong cùng một giao dịch cục bộ, rồi báo thành công; (5) nếu cổng không phản hồi hoặc báo lỗi, lượt quét được đánh dấu mở cổng thất bại và kiosk không báo thành công.
    \item Khi Gate Agent khởi động lại, các lượt quét còn ở trạng thái chờ xác nhận được đưa lên bảng điều khiển của bảo vệ để đối soát.
    \item G4 gửi các bản ghi outbox theo lô lên C1; C1 lưu idempotent theo mã sự kiện (INV-06); G4 chỉ đánh dấu đã đồng bộ sau khi nhận xác nhận từ C1.
\end{itemize}

\noindent\textbf{3. Hệ quả.}
\begin{itemize}
    \item \textit{Lợi ích:} không mất dấu vết kể cả khi mất nguồn giữa chừng; không có giao dịch chấm công cho lượt cổng không mở; dữ liệu được chuyển giao ít nhất một lần và không trùng lặp.
    \item \textit{Đánh đổi:} mỗi lượt quét có hai lần ghi đĩa đồng bộ (khoảng 50 ms mỗi lần, nằm trong ngân sách ở Bảng~\ref{tab:time-budget}); máy chủ biên cần SSD công nghiệp; cần cơ chế đối soát lượt quét chờ xác nhận.
\end{itemize}

\noindent\textbf{4. Phương án thay thế.}
\begin{itemize}
    \item \textit{Mở cổng trước, ghi dữ liệu sau vào bộ nhớ:} bị loại vì mất dữ liệu khi mất nguồn (vi phạm NFR-REL-03).
    \item \textit{Ghi giao dịch trước khi mở cổng:} bị loại vì tạo giao dịch cho lượt cổng không mở được và phải báo thành công trước khi cổng xác nhận (vi phạm NFR-AVL-03, R-05).
    \item \textit{Giao dịch phân tán (two-phase commit) giữa cổng và trung tâm:} bị loại vì không hoạt động khi mất kết nối.
\end{itemize}

%------------------------------------------------------------------------------
\subsubsection{ADR-04: Architecture style kết hợp}\label{adr:04}
\textbf{Trạng thái:} Accepted. \textbf{Nhóm:} cấu trúc.

\noindent\textbf{1. Bối cảnh.} Phân tích quantum (Mục~\ref{sec:quantum}) cho thấy hệ thống cần hai loại đơn vị triển khai độc lập. Việc đánh giá các style theo tiêu chí K1--K8 (Mục~\ref{sec:style-evaluation}) cho thấy không style đơn lẻ nào đáp ứng tốt đồng thời yêu cầu hoạt động khi mất kết nối (K4) và chi phí, năng lực của nhóm (K6, K7).

\noindent\textbf{2. Quyết định.}
\begin{itemize}
    \item Central Quantum là \textbf{modular monolith}: một đơn vị triển khai gồm các module C1--C9, mỗi module sở hữu dữ liệu của mình và chỉ giao tiếp với module khác qua interface công khai hoặc sự kiện trong tiến trình.
    \item Gate Agent áp dụng \textbf{microkernel kết hợp ports and adapters}: lõi xử lý xác thực và lưu trữ; thiết bị quét và cổng xoay của từng nhà cung cấp được cắm vào qua adapter cài đặt interface do lõi định nghĩa.
    \item Hai quantum liên kết bằng transactional outbox (ADR-03, ADR-05), không dùng message broker.
\end{itemize}

\noindent\textbf{3. Hệ quả.}
\begin{itemize}
    \item \textit{Lợi ích:} giữ giao dịch ACID cho các bất biến tại trung tâm (INV-01, INV-08, INV-09); chi phí vận hành thấp vì chỉ có hai loại đơn vị triển khai; thay nhà cung cấp cổng xoay chỉ cần thêm adapter (NFR-MOD-01).
    \item \textit{Đánh đổi:} cần kỷ luật giữ ranh giới module, được kiểm tra tự động (ADR-07); nâng cấp một module tại trung tâm phải triển khai lại toàn bộ Central, điều này chấp nhận được vì cổng vẫn hoạt động độc lập.
\end{itemize}

\noindent\textbf{4. Phương án thay thế.} Layered (55 điểm), modular monolith đơn thuần (65), microkernel (59), service-based với 4--6 dịch vụ (72), event-driven có broker (63) và microservices (54) đều được đánh giá thấp hơn phương án kết hợp (85) theo Bảng~\ref{tab:style-scores}.

%------------------------------------------------------------------------------
\subsubsection{ADR-05: Giao tiếp giữa hai quantum và bên trong Central}\label{adr:05}
\textbf{Trạng thái:} Accepted. \textbf{Nhóm:} giao tiếp.

\noindent\textbf{1. Bối cảnh.} Hệ thống có hai loại ranh giới giao tiếp: giữa Gate Quantum và Central Quantum qua WAN, và giữa các module bên trong Central. Giao tiếp giữa hai quantum không được làm chậm luồng quét và phải chịu được mất kết nối; giao tiếp bên trong Central phải giữ ghép nối lỏng và không có vòng phụ thuộc.

\noindent\textbf{2. Quyết định.}
\begin{itemize}
    \item Giữa hai quantum chỉ dùng hợp đồng dữ liệu \textit{sync-contract} có quản lý phiên bản, truyền bằng \textbf{REST/JSON qua HTTPS}. Chiều lên: G4 gửi lô sự kiện từ outbox, idempotent theo mã sự kiện. Chiều xuống: G4 định kỳ kéo phần thay đổi của credential, khóa công khai và cấu hình trạm, bảo đảm độ trễ $\le$ 60 giây (NFR-SEC-06); thay đổi thu hồi credential được ưu tiên.
    \item Bên trong Central: truy vấn qua interface công khai của module; thông báo thay đổi trạng thái qua sự kiện trong tiến trình. Chiều phụ thuộc giữa các module phải không có vòng (Mục~\ref{sec:module-views}); để giữ quy tắc này, C1 gọi interface của C3 và C5 gọi interface của C2 thay vì phát sự kiện ngược chiều.
\end{itemize}

\noindent\textbf{3. Hệ quả.}
\begin{itemize}
    \item \textit{Lợi ích:} Central có thể bảo trì mà không dừng cổng; module phát sự kiện không cần biết module nhận; REST/JSON dễ kiểm thử và gỡ lỗi.
    \item \textit{Đánh đổi:} dữ liệu giữa các quantum và read model báo cáo chỉ nhất quán cuối cùng; phải quản lý phiên bản \textit{sync-contract} để Gate Agent và Central khác phiên bản vẫn tương thích trong thời gian nâng cấp.
\end{itemize}

\noindent\textbf{4. Phương án thay thế.}
\begin{itemize}
    \item \textit{Message broker (Kafka, RabbitMQ):} bị loại vì quá mức cần thiết với tải vài sự kiện mỗi giây và thêm thành phần cần vận hành, bảo mật.
    \item \textit{Chia sẻ CSDL giữa hai quantum:} bị loại vì không hoạt động khi mất kết nối và phá ranh giới quantum.
    \item \textit{gRPC:} không chọn vì REST/JSON đủ đáp ứng tải hiện tại và đơn giản hơn cho việc kiểm thử, giám sát.
\end{itemize}

%------------------------------------------------------------------------------
\subsubsection{ADR-06: Tích hợp cổng xoay và HR/Payroll qua adapter}\label{adr:06}
\textbf{Trạng thái:} Accepted. \textbf{Nhóm:} tích hợp, chịu lỗi.

\noindent\textbf{1. Bối cảnh.} Hệ thống phải tích hợp với hệ thống cổng xoay và hệ thống HR/Payroll qua API (C-01), cả hai nằm ngoài tầm kiểm soát của nhóm. Cổng xoay có thể không phản hồi hoặc báo lỗi; HR/Payroll có thể tạm thời không khả dụng đúng lúc truyền dữ liệu cuối tháng (R-03, C-02).

\noindent\textbf{2. Quyết định.}
\begin{itemize}
    \item \textbf{Cổng xoay:} domain của G1 định nghĩa port \texttt{TurnstileGateway}; mỗi nhà cung cấp có một adapter trong G2. Lời gọi mở cổng có giới hạn thời gian chờ 800 ms, lớn hơn thời gian phản hồi cam kết 500 ms (A-04). Khi hết thời gian chờ, G2 kiểm tra trạng thái cổng trước khi gửi lại lệnh (NFR-REL-02); không gửi lại theo backoff để tránh mở cổng hai lần. Sau năm lỗi liên tiếp, circuit breaker ngừng gửi lệnh tự động và cảnh báo bảo vệ để mở cổng thủ công (FR-10).
    \item \textbf{HR/Payroll:} C5 cài đặt hai port \texttt{HrPayroll\-Gateway} (gửi đợt dữ liệu chấm công) và \texttt{Employee\-Source\-Gateway} (nhận dữ liệu nhân viên, ADR-09). Đợt truyền được gửi bất đồng bộ, gửi lại theo exponential backoff (1, 2, 4, 8 phút\ldots) khi lỗi tạm thời (NFR-REL-01), idempotent theo mã đợt (NFR-INT-04). Nhân viên HR có thể theo dõi và gửi lại thủ công (FR-22); hệ thống cảnh báo khi gần hạn 12:00 ngày 1 mà đợt chính chưa được xác nhận.
\end{itemize}

\noindent\textbf{3. Hệ quả.}
\begin{itemize}
    \item \textit{Lợi ích:} sự cố của hệ thống bên ngoài không làm treo luồng quét hay nghiệp vụ chốt kỳ; thay nhà cung cấp hoặc API bên ngoài chỉ cần thay adapter (NFR-MOD-01); adapter có thể thay bằng bản giả lập khi kiểm thử.
    \item \textit{Đánh đổi:} khi circuit breaker ngắt, việc ra vào tại trạm đó phụ thuộc thao tác thủ công của bảo vệ; cần màn hình theo dõi trạng thái truyền và cơ chế cảnh báo.
\end{itemize}

\noindent\textbf{4. Phương án thay thế.}
\begin{itemize}
    \item \textit{Gọi trực tiếp API bên ngoài trong mã nghiệp vụ:} bị loại vì làm nghiệp vụ phụ thuộc vào API bên ngoài (vi phạm NFR-MOD-01).
    \item \textit{Ghi trực tiếp vào CSDL của HR/Payroll:} bị loại vì vi phạm ranh giới hệ thống và có rủi ro về toàn vẹn dữ liệu của hệ thống HR/Payroll.
\end{itemize}

%------------------------------------------------------------------------------
\subsubsection{ADR-07: Tổ chức mã nguồn và kiểm thử kiến trúc}\label{adr:07}
\textbf{Trạng thái:} Accepted. \textbf{Nhóm:} phát triển, khả năng sửa đổi.

\noindent\textbf{1. Bối cảnh.} Hệ thống gồm hai ứng dụng triển khai riêng (\texttt{gate-agent}, \texttt{central}) dùng chung hợp đồng \textit{sync-contract}. Nhóm phát triển nhỏ cần đồng bộ thay đổi trên hợp đồng này một cách dễ dàng, đồng thời cấu trúc mã nguồn phải phản ánh ranh giới domain để tránh suy thoái kiến trúc theo thời gian.

\noindent\textbf{2. Quyết định.}
\begin{itemize}
    \item Toàn bộ mã nguồn nằm trong một kho mã (monorepo) \texttt{tmc}, gồm \texttt{sync-contract}, \texttt{gate-agent}, \texttt{central}, mã nguồn giao diện \texttt{web} và cấu hình triển khai \texttt{deploy} (Bảng~\ref{tab:implementation-view}).
    \item Mã nguồn được chia package theo module (package-by-feature), mỗi module ứng với một component; bên trong mỗi module có bốn tầng \texttt{api}, \texttt{application}, \texttt{domain}, \texttt{infrastructure} (Mục~\ref{sec:module-views}).
    \item Các quy tắc phụ thuộc ở Mục~\ref{sec:module-views} được kiểm tra tự động bằng ArchUnit trong quy trình tích hợp liên tục (CI); vi phạm làm build thất bại.
    \item Tham số thay đổi theo nhà máy hoặc môi trường (ngưỡng mất kết nối, giới hạn thời gian chờ, chu kỳ đồng bộ) được nạp từ cấu hình (C9 hoặc biến môi trường), không viết cứng trong mã.
\end{itemize}

\noindent\textbf{3. Hệ quả.}
\begin{itemize}
    \item \textit{Lợi ích:} một thay đổi có thể cập nhật đồng thời \textit{sync-contract}, phía gửi và phía nhận; các thành viên làm việc song song trên các module khác nhau; ranh giới kiến trúc được giữ bằng kiểm thử tự động thay vì chỉ dựa vào kỷ luật.
    \item \textit{Đánh đổi:} kho mã lớn dần; CI cần được cấu hình để chỉ build và kiểm thử phần bị thay đổi.
\end{itemize}

\noindent\textbf{4. Phương án thay thế.}
\begin{itemize}
    \item \textit{Mỗi ứng dụng một kho mã:} bị loại vì việc quản lý phiên bản hợp đồng dùng chung phức tạp với nhóm nhỏ.
    \item \textit{Chia package theo tầng kỹ thuật (controller, service, repository):} bị loại vì phá vỡ ranh giới domain và không khớp với cách phân rã đã chọn (Mục~\ref{sec:partitioning}).
\end{itemize}

%------------------------------------------------------------------------------
\subsubsection{ADR-08: Kiểm soát truy cập, nhật ký kiểm toán bất biến và bảo vệ dữ liệu cá nhân}\label{adr:08}
\textbf{Trạng thái:} Accepted. \textbf{Nhóm:} bảo mật, quyền riêng tư.

\noindent\textbf{1. Bối cảnh.} Hệ thống xử lý dữ liệu cá nhân của nhân viên (họ tên, mã nhân viên, bộ phận, lịch sử ra vào) và phải tuân thủ quy định về bảo vệ dữ liệu cá nhân (C-03, Nghị định 13/2023/NĐ-CP). Mọi chức năng được bảo vệ phải kiểm tra quyền (NFR-SEC-01), mọi thao tác quan trọng phải để lại nhật ký không thể sửa xóa (NFR-AUD-01 -- NFR-AUD-05), và thiết bị đặt tại nhà máy có thể bị can thiệp.

\noindent\textbf{2. Quyết định.}
\begin{itemize}
    \item \textbf{Kiểm soát truy cập:} C7 quản lý tài khoản, vai trò và quyền. Mọi API của Central yêu cầu token có chữ ký (JWT, thuật toán RS256) và kiểm tra quyền ở cả mức chức năng và mức dữ liệu; giám sát sản xuất chỉ xem được nhân viên thuộc bộ phận mình phụ trách (FR-19, INV-11).
    \item \textbf{Thiết bị:} kiosk không yêu cầu tài khoản người dùng. Mỗi Gate Agent có chứng chỉ thiết bị riêng để xác thực hai chiều (mTLS) với trung tâm; chứng chỉ của thiết bị bị mất được thu hồi.
    \item \textbf{Bảo vệ đăng nhập khi mất kết nối:} thông tin xác thực (dạng băm) và quyền mở cổng thủ công của các tài khoản bảo vệ được đồng bộ xuống Gate Agent; thao tác mở cổng thủ công được ghi nhật ký cục bộ và đồng bộ lên C8.
    \item \textbf{Bảo vệ dữ liệu:} bản sao tại Gate Agent chỉ chứa dữ liệu cần cho xác thực (mã credential, mã nhân viên, trạng thái); CSDL cục bộ được mã hóa; CSDL trung tâm được mã hóa khi lưu. Kết nối qua WAN và tới trung tâm dùng TLS $\ge$ 1.2, ưu tiên 1.3; kết nối tới cổng xoay theo hợp đồng của nhà cung cấp trong VLAN cô lập (NFR-SEC-03). Khóa bí mật ký mã QR chỉ nằm trong kho khóa của trung tâm (ADR-02).
    \item \textbf{Nhật ký kiểm toán:} C8 chỉ cho phép thêm bản ghi, không tài khoản ứng dụng nào có quyền sửa hoặc xóa; mỗi bản ghi lưu giá trị băm của bản ghi trước để phát hiện can thiệp trực tiếp vào CSDL; dữ liệu được lưu trực tuyến tối thiểu 12 tháng trước khi chuyển sang kho lưu trữ dài hạn (INV-10).
\end{itemize}

\noindent\textbf{3. Hệ quả.}
\begin{itemize}
    \item \textit{Lợi ích:} đáp ứng yêu cầu pháp lý và yêu cầu kiểm toán; thiết bị tại nhà máy bị đánh cắp không làm lộ dữ liệu và không thể dùng để gửi dữ liệu giả hay tạo mã QR; nhật ký kiểm toán là bằng chứng tin cậy khi giải quyết khiếu nại.
    \item \textit{Đánh đổi:} mã hóa làm tăng nhẹ thời gian đọc ghi; cần quy trình quản lý chứng chỉ và khóa; bảng nhật ký lớn dần theo thời gian, cần phân vùng theo tháng.
\end{itemize}

\noindent\textbf{4. Phương án thay thế.}
\begin{itemize}
    \item \textit{Ghi nhật ký vào tệp văn bản:} bị loại vì dễ bị sửa hoặc xóa bởi người có quyền quản trị máy chủ.
    \item \textit{Lưu đầy đủ hồ sơ nhân viên tại kiosk để hiển thị:} bị loại vì làm tăng rủi ro lộ dữ liệu cá nhân tại hiện trường.
\end{itemize}

%------------------------------------------------------------------------------
\subsubsection{ADR-09: Đồng bộ dữ liệu nhân viên từ HR/Payroll}\label{adr:09}
\textbf{Trạng thái:} Accepted. \textbf{Nhóm:} tích hợp, toàn vẹn dữ liệu.

\noindent\textbf{1. Bối cảnh.} HR/Payroll là nguồn dữ liệu nhân sự chính thức; hệ thống TMC chỉ tiếp nhận và đồng bộ, không tự tạo dữ liệu nhân sự gốc (Mục~\ref{sec:external-boundaries}, FR-01). Khi nhân viên nghỉ việc, QR credential phải bị thu hồi kịp thời (NFR-SEC-06). Khả năng API của HR/Payroll cần được xác nhận (A-03).

\noindent\textbf{2. Quyết định.}
\begin{itemize}
    \item C5 định kỳ (mặc định 15 phút, cấu hình tại C9) lấy phần thay đổi dữ liệu nhân viên qua port \texttt{Employee\-Source\-Gateway}, rồi gọi interface của C2 để cập nhật hồ sơ theo mã nhân viên; việc cập nhật là idempotent nên chạy lại không tạo dữ liệu trùng.
    \item Khi nhân viên chuyển sang trạng thái không hoạt động, C2 thu hồi QR credential đang hiệu lực và thay đổi được phát xuống các trạm theo luồng đồng bộ chiều xuống.
    \item Khi API của HR/Payroll chưa sẵn sàng, nhân viên HR có thể cập nhật thủ công; mọi cập nhật được ghi nhật ký kiểm toán.
\end{itemize}

\noindent\textbf{3. Hệ quả.}
\begin{itemize}
    \item \textit{Lợi ích:} một nguồn dữ liệu nhân sự duy nhất, tránh lệch dữ liệu giữa hai hệ thống; việc thu hồi credential khi nhân viên nghỉ việc diễn ra tự động.
    \item \textit{Đánh đổi:} có độ trễ tối đa bằng chu kỳ đồng bộ giữa HR/Payroll và TMC; phụ thuộc vào khả năng cung cấp dữ liệu thay đổi của API HR/Payroll.
\end{itemize}

\noindent\textbf{4. Phương án thay thế.}
\begin{itemize}
    \item \textit{Nhập dữ liệu thủ công hoặc qua tệp CSV:} bị loại làm phương án chính vì chậm và dễ sai; chỉ dùng làm phương án dự phòng.
    \item \textit{HR/Payroll chủ động đẩy thay đổi sang TMC (webhook):} không chọn vì phụ thuộc vào khả năng của hệ thống hiện có, chưa được xác nhận (A-03).
    \item \textit{Đọc trực tiếp CSDL của HR/Payroll:} bị loại vì vi phạm ranh giới hệ thống.
\end{itemize}

%------------------------------------------------------------------------------
\subsubsection{ADR-10: Nền tảng công nghệ}\label{adr:10}
\textbf{Trạng thái:} Accepted. \textbf{Nhóm:} công nghệ.

\noindent\textbf{1. Bối cảnh.} Kiến trúc đã chọn cần một nền tảng hỗ trợ tốt việc tổ chức module, interface và kiểm thử kiến trúc (ADR-04, ADR-07), chạy được trên máy chủ biên và máy chủ trung tâm, phù hợp với năng lực của nhóm (K7) và có chi phí vận hành thấp (K6).

\noindent\textbf{2. Quyết định.}
\begin{itemize}
    \item \textbf{Backend:} Java 21 và Spring Boot 3 cho cả \texttt{gate-agent} và \texttt{central}; \textit{sync-contract} là thư viện Java dùng chung. Central dùng Spring Modulith để kiểm tra ranh giới module, phát sự kiện trong tiến trình và lưu sự kiện chưa xử lý (dùng làm outbox trung tâm).
    \item \textbf{Giao diện:} React cho giao diện kiosk, bảng điều khiển của bảo vệ và giao diện quản trị; bản build tĩnh được đóng gói vào artifact \texttt{gate-agent} hoặc \texttt{central} tương ứng.
    \item \textbf{Lưu trữ:} SQLite được mã hóa tại Gate Agent, truy cập qua JDBC; PostgreSQL tại Central, truy cập qua Spring Data JPA.
    \item \textbf{Hạ tầng và thư viện khác:} container cho cả hai loại node; Nginx làm reverse proxy; thuật toán Ed25519 có sẵn trong JDK để ký và kiểm tra mã QR; ArchUnit cho kiểm thử kiến trúc.
\end{itemize}

\noindent\textbf{3. Hệ quả.}
\begin{itemize}
    \item \textit{Lợi ích:} một ngôn ngữ cho toàn bộ backend; hệ sinh thái Spring hỗ trợ sẵn tiêm phụ thuộc, giao dịch, sự kiện và kiểm thử, phù hợp với ports and adapters; Spring Modulith và ArchUnit giúp giữ ranh giới kiến trúc; SQLite không cần cài đặt máy chủ CSDL tại nhà máy.
    \item \textit{Đánh đổi:} ứng dụng Java tiêu tốn bộ nhớ hơn và khởi động chậm hơn, chấp nhận được với cấu hình máy chủ biên ở Bảng~\ref{tab:nodes}; nhóm phải làm việc với hai ngôn ngữ (Java cho backend, JavaScript/TypeScript cho React).
\end{itemize}

\noindent\textbf{4. Phương án thay thế.}
\begin{itemize}
    \item \textit{TypeScript với NestJS cho backend:} không chọn; ưu điểm là dùng một ngôn ngữ cho cả backend và giao diện, nhưng công cụ kiểm thử kiến trúc và quản lý module kém hơn hệ sinh thái Java, và không phù hợp năng lực hiện có của nhóm.
    \item \textit{Python với FastAPI:} bị loại vì kiểu động khiến việc thể hiện và kiểm tra interface, ranh giới module khó hơn.
    \item \textit{CSDL máy chủ (PostgreSQL) tại nhà máy:} không chọn vì tăng chi phí vận hành tại hiện trường trong khi tải tại nhà máy thấp.
\end{itemize}
